\scanpage{32}
% archive foliation: 31

\margin{\uline{NB}
\[
\begin{array}{lll}
\overset{C_Q}{\underset{\simeq}{\underline\omega}}\subset\mathbb{F}_2^{\underline\omega}=\mathbb{F}_2^C=V_S\subset\mathfrak D_Q
&\underline\omega\simeq V_S\setminus Z_Q&\longleftrightarrow\mathrm{sg}(g)\\
D_Q\simeq\underline\omega\wedge\underline\mu\subset\mathbb{F}_2^D=V_A\subset\mathfrak D_Q
&\underline\omega\wedge\underline\mu\simeq V_A\setminus Z_Q&\longleftrightarrow\mathrm{sg}(g)\det(g)\\
\underline\omega_Q\simeq\underline\mu\subset\mathfrak D_Q^+\subset\mathfrak D_Q
&\underline\mu\simeq\mathfrak D_Q^+\setminus Z_Q&\longleftrightarrow\det(g)
\end{array}
\]
où $\underline\mu\unsure{\simeq|=}(\bigwedge^2_\Lambda M)^*=\mathrm{Or}(M)$,
$Z_Q=\{1,\underline a_Q\}=\mathrm{Cent}\,\mathfrak D_Q$.}

\margin{$GL_\Lambda(M)'\simeq\unsure{\mathcal{G}_I|\mathfrak S_I}\times\mathbb{F}_2
\simeq\mathrm{Aut}(\Delta)$,
où $\Delta$ est l'octaèdre non orienté associé à $M$, de sommets $\Delta=M^*/\{\pm1\}\ldots$}

\[
(28)\qquad
\diagram{\begin{tikzcd}[column sep=small,row sep=large]
&&&GL_\Lambda(M)[\simeq\unsure{\mathcal{G}_I|\mathfrak S_I}\cdot\mathbb{F}_2^{\underline\omega}]&\\
&&\unsure{\mathcal{G}_I^+|\mathfrak S_I^+}\times\mathbb{F}_2^{\underline\omega}\arrow[ur,"2"]&
[\unsure{\mathcal{G}_I|\mathfrak S_I}\times Z]\arrow[u,dashed]&SL_\Lambda(M)\arrow[ul,"2"']\\
&M_0\times\mathbb{F}_2^{\underline\omega}\simeq\mathfrak{gl}(M_0)\arrow[ur,"3"]&&
\unsure{\mathcal{G}_I^+|\mathfrak S_I^+}\times Z\arrow[ul,"2"]\arrow[u,dashed]\arrow[ur,"2"']&\\
&&M_0\times Z\simeq\mathfrak{sl}(M_0)\arrow[ul,"2"']\arrow[ur,"3"]&&
\unsure{\mathcal{G}_I^+|\mathfrak S_I^+}\arrow[ul,"2"']\\
\mathbb{F}_2^{\underline\omega}\arrow[uur,"4"]&&&M_0\arrow[ul,"2"']\arrow[ur,"3"]&\\
&Z\arrow[ul,"2"]\arrow[uur,"4"]&&&\\
&&1\arrow[ul,"2"]\arrow[uur,"4"']&&
\arrow[from=1-4,to=4-5,-,decorate,decoration={brace,amplitude=6pt},bend left=35,"\mathfrak D_Q"]
\arrow[from=2-5,to=4-5,-,decorate,decoration={brace,amplitude=5pt},bend left=18,
"{\substack{\mathfrak D_Q^+\simeq(SL_\Lambda(M))_{\mathrm{ab}}\\
\simeq\bigwedge^2_\Lambda M\\\deleted{\unsure{\mathbb{F}_2}}}}"]
\arrow[from=4-5,to=5-4,-,decorate,decoration={brace,amplitude=5pt},bend left=25,
"{\unsure{\mathcal{G}_J^+|\mathfrak S_J^+}}"]
\arrow[from=5-1,to=1-4,-,decorate,decoration={brace,amplitude=6pt},bend left=30,
"{\unsure{\mathcal{G}_I|\mathfrak S_I}}"']
\arrow[from=3-2,to=1-4,-,decorate,decoration={brace,amplitude=5pt},bend left=30,
"{\substack{\unsure{\mathcal{G}_J|\mathfrak S_J}\\\simeq GL_{\mathbb{F}_2}(M_0)}}"']
\end{tikzcd}}
\]

\uline{NB} Dans ce diagramme, toutes les flèches sont injections, et les cinq carrés qui y apparaissent, donc aussi les composés, sont cartésiens. On a marqué sur chaque flèche l'indice \deleted{\unsure{pour}} l'inclusion de sous-groupes qu'elle représente.

Les iso.
\[
GL_\Lambda(M)\simeq\unsure{\mathcal{G}_I|\mathfrak S_I}\cdot\mathbb{F}_2^{\underline\omega}\quad(1/2\text{ direct})\quad\text{et}
\]
\[
\mathrm{Ker}(\mathrm{sg}\cdot\det)\simeq\unsure{\mathcal{G}_I|\mathfrak S_I}\times\mathbb{F}_2\quad(\text{prod direct}
\]
sont valables quand $M$ est ``$1/2$ déployé'' i.e. le torseur $P_M$ est trivialisé, alors on a
\[
GL_\Lambda(M)\simeq
\underset{\unsure{\mathcal{G}_J|\mathfrak S_J}}{\underbrace{GL_{\mathbb{F}_2}(M_0)}}\cdot(M_0\times\mathbb{F}_2^{\underline\omega})
\simeq
\underset{\unsure{\mathcal{G}_I|\mathfrak S_I}}{\underbrace{(GL_{\mathbb{F}_2}(M_0)\cdot M_0)}}\cdot\mathbb{F}_2^{\underline\omega}
\]
mais ces isom. dépendent du choix d'une trivialisation de $P_M\ldots$
